\section{Números complejos}

\begin{teoria}
\textbf{Unidad imaginaria y forma binómica.}
La unidad imaginaria $i$ cumple $i^2 = -1$ (y $(-i)^2 = -1$). Un número complejo se escribe en \emph{forma binómica} o rectangular como
\[
z = a + i\,b, \qquad a,b \in \R,
\]
donde $a = \operatorname{Re}(z)$ es la parte real y $b = \operatorname{Im}(z)$ la parte imaginaria. El conjunto es
\[
\C = \{ a + i\,b : a,b \in \R \} = \R + i\,\R, \qquad \R \subset \C.
\]
Potencias de $i$: son cíclicas de período $4$,
\[
i^0 = 1, \quad i^1 = i, \quad i^2 = -1, \quad i^3 = -i, \quad i^4 = 1,
\]
y en general $i^{2n} = (-1)^n$, $\; i^{2n+1} = (-1)^n\, i$ para todo $n \in \N$.
\end{teoria}

\begin{teoria}
\textbf{Operaciones en forma binómica.} Sean $z = a_z + i\,b_z$ y $w = a_w + i\,b_w$.
\begin{align*}
z = w &\iff a_z = a_w \ \wedge\ b_z = b_w, \\
z + w &= (a_z + a_w) + i\,(b_z + b_w), \\
z\,w &= (a_z a_w - b_z b_w) + i\,(a_z b_w + a_w b_z).
\end{align*}
El producto se obtiene distribuyendo y usando $i^2 = -1$. El inverso multiplicativo de $z \neq 0$ y la división son
\[
z^{-1} = \frac{a_z}{a_z^2 + b_z^2} + i\,\frac{-b_z}{a_z^2 + b_z^2}, \qquad
\frac{z}{w} = z\,w^{-1}.
\]
En la práctica, para dividir se multiplica y divide por el conjugado del denominador:
\[
\frac{z}{w} = \frac{z\,\bar w}{w\,\bar w} = \frac{z\,\bar w}{\abs{w}^2}.
\]
\textbf{Observación.} $\C$ es un cuerpo \emph{no ordenado}: no existe un orden en $\C$ que respete los axiomas de orden de $\R$ (tricotomía, transitividad, monotonía).
\end{teoria}

\begin{teoria}
\textbf{Módulo, conjugado y argumento.} Sea $z = a + i\,b$.
\[
\abs{z} = \sqrt{a^2 + b^2} \quad (\text{módulo, } \rho = \abs{z}), \qquad
\bar z = a - i\,b \quad (\text{conjugado}).
\]
El \emph{argumento} $\theta = \arg(z)$ es el ángulo que forma $z$ con el semieje real positivo, $0 \le \theta < 2\pi$. Se calcula con $\arctan(b/a)$ ajustando por cuadrante:
\[
\arg(z) = \begin{cases}
\arctan\!\left(\dfrac{b}{a}\right) & a > 0,\ b > 0 \quad (\text{I}), \\[2mm]
\pi - \arctan\!\left(\dfrac{\abs{b}}{\abs{a}}\right) & a < 0,\ b > 0 \quad (\text{II}), \\[2mm]
\pi + \arctan\!\left(\dfrac{\abs{b}}{\abs{a}}\right) & a < 0,\ b < 0 \quad (\text{III}), \\[2mm]
2\pi - \arctan\!\left(\dfrac{\abs{b}}{\abs{a}}\right) & a > 0,\ b < 0 \quad (\text{IV}),
\end{cases}
\]
y en los ejes: $\arg = 0$ si $b=0,\,a>0$; $\arg = \tfrac{\pi}{2}$ si $a=0,\,b>0$; $\arg = \pi$ si $b=0,\,a<0$; $\arg = \tfrac{3\pi}{2}$ si $a=0,\,b<0$.

\medskip
\textbf{Propiedades del módulo.}
\[
\abs{z} \ge 0,\quad \abs{z}=0 \iff z=0; \qquad \abs{z\,w} = \abs{z}\,\abs{w}; \qquad \abs{\tfrac{z}{w}} = \frac{\abs{z}}{\abs{w}};
\]
\[
\abs{z+w} \le \abs{z}+\abs{w}; \qquad \abs{\operatorname{Re}(z)} \le \abs{z},\quad \abs{\operatorname{Im}(z)} \le \abs{z}.
\]
\textbf{Propiedades del conjugado.}
\[
z\,\bar z = \abs{z}^2; \qquad z + \bar z = 2\operatorname{Re}(z); \qquad z - \bar z = 2i\operatorname{Im}(z);
\]
\[
\overline{z+w} = \bar z + \bar w; \qquad \overline{z\,w} = \bar z\,\bar w; \qquad \abs{\bar z} = \abs{z}; \qquad \bar z = z \iff \operatorname{Im}(z)=0.
\]
En particular $\operatorname{Re}(z) = \dfrac{z + \bar z}{2}$ y $\operatorname{Im}(z) = \dfrac{z - \bar z}{2i}$.
\end{teoria}

\begin{teoria}
\textbf{Forma polar, exponencial y fórmula de Euler.} Con $\rho = \abs{z}$ y $\theta = \arg(z)$:
\[
z = \rho\,(\cos\theta + i\,\operatorname{sen}\theta) \quad (\text{trigonométrica}), \qquad
z = \rho\, e^{i\theta} \quad (\text{exponencial}),
\]
donde la \emph{fórmula de Euler} establece
\[
e^{i\theta} = \cos\theta + i\,\operatorname{sen}\theta \qquad \forall\,\theta \in \R.
\]
De ahí se obtienen $\cos\theta = \dfrac{e^{i\theta} + e^{-i\theta}}{2}$ y $\operatorname{sen}\theta = \dfrac{e^{i\theta} - e^{-i\theta}}{2i}$.

Igualdad en forma exponencial: $z = w \iff \abs{z}=\abs{w} \ \wedge\ \arg(z) = \arg(w) + 2k\pi,\ k \in \Z$.

\medskip
\textbf{Producto, cociente y argumento.} La forma exponencial simplifica el producto:
\[
z\,w = \abs{z}\,\abs{w}\,e^{\,i(\arg z + \arg w)}, \qquad
\frac{z}{w} = \frac{\abs{z}}{\abs{w}}\,e^{\,i(\arg z - \arg w)}.
\]
\[
\arg(z\,w) = \arg z + \arg w, \quad \arg\!\left(\tfrac{z}{w}\right) = \arg z - \arg w, \quad \arg(z^n) = n\arg z, \quad \arg(z^{-1}) = -\arg z.
\]
\end{teoria}

\begin{teoria}
\textbf{Potenciación y fórmula de De Moivre.} Para $z = \rho\,e^{i\theta}$ y $n \in \N$:
\[
z^n = \rho^n\, e^{\,i\,n\theta} = \rho^n\big(\cos(n\theta) + i\,\operatorname{sen}(n\theta)\big).
\]
En particular, para $\rho = 1$ se obtiene la \emph{fórmula de De Moivre}:
\[
(\cos\theta + i\,\operatorname{sen}\theta)^n = \cos(n\theta) + i\,\operatorname{sen}(n\theta).
\]
\end{teoria}

\begin{teoria}
\textbf{Radicación: raíces $n$-ésimas.} La ecuación $w^n = z$, con $z = \rho\,e^{i\theta}$ y $n \in \N$, tiene exactamente $n$ soluciones distintas:
\[
w_k = \sqrt[n]{\rho}\; e^{\,i\left(\frac{\theta + 2k\pi}{n}\right)}, \qquad k = 0, 1, \dots, n-1.
\]
Las $n$ raíces se ubican en los vértices de un polígono regular de $n$ lados sobre la circunferencia de radio $\sqrt[n]{\rho}$; se reparten con paso angular $\tfrac{2\pi}{n}$ y para $k \ge n$ se repiten cíclicamente ($w_n = w_0$, etc.).

\medskip
\textbf{Caso raíz cuadrada ($n=2$).} Las soluciones de $w^2 = z$ cumplen
\[
\operatorname{Re}(w) = \pm\sqrt{\frac{\abs{z} + \operatorname{Re}(z)}{2}}, \qquad
\operatorname{Im}(w) = \pm\sqrt{\frac{\abs{z} - \operatorname{Re}(z)}{2}},
\]
eligiendo los signos de modo que $2\operatorname{Re}(w)\operatorname{Im}(w) = \operatorname{Im}(z)$.
\end{teoria}

\begin{teoria}
\textbf{Logaritmo complejo y potencia compleja.} La ecuación $e^w = z$ (con $z \neq 0$) tiene infinitas soluciones:
\[
w_k = \ln\abs{z} + i\,\big(\arg(z) + 2k\pi\big), \qquad k \in \Z.
\]
El \emph{valor principal} ($\operatorname{vp}$) toma $k=0$:
\[
\operatorname{vp}\,\ln z = \ln\abs{z} + i\,\arg(z).
\]
La \emph{potencia compleja} se define vía el logaritmo:
\[
z^w = e^{\,w\,\ln z}.
\]
Como $\ln z$ es multivaluado, $z^w$ lo es en general; el valor principal usa $\operatorname{vp}\,\ln z$.
\end{teoria}

\begin{receta}
\textbf{Hallar las $n$ raíces $n$-ésimas de $z$ (resolver $w^n = z$).}
\begin{enumerate}
  \item Pasar $z$ a forma polar: calcular $\rho = \abs{z} = \sqrt{a^2+b^2}$ y $\theta = \arg(z)$ (¡cuidar el cuadrante!).
  \item El módulo de todas las raíces es $\sqrt[n]{\rho}$ (raíz real positiva).
  \item Los argumentos son $\dfrac{\theta + 2k\pi}{n}$ para $k = 0, 1, \dots, n-1$ (exactamente $n$ valores).
  \item Escribir $w_k = \sqrt[n]{\rho}\; e^{\,i\left(\frac{\theta + 2k\pi}{n}\right)}$ y, si se pide, pasar cada $w_k$ a binómica con Euler.
  \item Verificación: las raíces están igualmente espaciadas ($\tfrac{2\pi}{n}$) sobre la circunferencia de radio $\sqrt[n]{\rho}$.
\end{enumerate}
\end{receta}

\begin{receta}
\textbf{Resolver la ecuación exponencial $e^w = z$ (logaritmo complejo).}
\begin{enumerate}
  \item Escribir $w = \alpha + i\beta$ y pasar $z$ a polar: $z = \rho\,e^{i\theta}$.
  \item Como $e^w = e^\alpha e^{i\beta}$, igualar módulos y argumentos: $e^\alpha = \rho$ y $\beta = \theta + 2k\pi$.
  \item Despejar: $\alpha = \ln\rho = \ln\abs{z}$ y $\beta = \theta + 2k\pi$, con $k \in \Z$.
  \item Escribir la familia $w_k = \ln\abs{z} + i(\theta + 2k\pi)$. El valor principal es $k=0$.
\end{enumerate}
\end{receta}

\begin{ejemplo}{Resueltos Complejos, Ej.\ 1 (Guía I)}
Sean $z_1 = i$, $z_2 = 1+i$, $z_4 = 81$, $z_5 = -21$. Calcular varias operaciones.

\smallskip
\textbf{Formas polares:} $z_1 = i$ tiene $\rho=1,\ \theta=\tfrac{\pi}{2}$; $\ z_2 = 1+i$ tiene $\rho=\sqrt2,\ \theta=\tfrac{\pi}{4}$; $\ z_4=81$ tiene $\rho=81,\ \theta=0$; $\ z_5=-21$ tiene $\rho=21,\ \theta=\pi$.

\smallskip
\textbf{(iii) Raíz cuarta de $z_4$:} $\sqrt[4]{z_4} = \sqrt[4]{81}\,e^{\,i\frac{0 + 2k\pi}{4}} = 3\,e^{\,i\frac{k\pi}{2}}$, $k=0,1,2,3$. En binómica:
\[
w_0 = 3, \quad w_1 = 3i, \quad w_2 = -3, \quad w_3 = -3i.
\]

\textbf{(v) $\operatorname{vp}\,\ln z_1$:} $\ln z_1 = \ln\abs{z_1} + i\theta_1 = \ln 1 + i\tfrac{\pi}{2} = i\tfrac{\pi}{2}$.

\smallskip
\textbf{(viii) $z_1^{245}$:} usando la ciclicidad de $i$,
\[
i^{245} = i^{244}\cdot i = (i^4)^{61}\cdot i = 1\cdot i = i.
\]

\textbf{(ix) $\operatorname{vp}\,z_1^{z_2}$:} por definición $z_1^{z_2} = e^{\,z_2 \ln z_1} = e^{\,(1+i)\cdot i\frac{\pi}{2}}$. Como $(1+i)\,i\tfrac{\pi}{2} = -\tfrac{\pi}{2} + i\tfrac{\pi}{2}$:
\[
z_1^{z_2} = e^{-\pi/2}\,e^{\,i\pi/2} = e^{-\pi/2}\big(\cos\tfrac{\pi}{2} + i\operatorname{sen}\tfrac{\pi}{2}\big) = i\,e^{-\pi/2}.
\]

\textbf{(x) $\operatorname{vp}\,z_4^{z_2}$:} $z_4^{z_2} = e^{\,(1+i)\ln 81} = e^{\ln 81}\,e^{\,i\ln 81} = 81\,e^{\,i\ln 81} = 81\cos(\ln 81) + 81 i\operatorname{sen}(\ln 81)$.
\end{ejemplo}

\begin{ejemplo}{Resueltos Complejos, Ej.\ 2}
Resolver $z^4 + 4iz^3 - 6z^2 - 4iz - 15 = 0$ en $\C$.

\smallskip
\textbf{Idea (binomio de Newton).} Reescribimos usando $i^2=-1,\ i^3=-i,\ i^4=1$:
\[
z^4 + 4iz^3 + 6i^2 z^2 + 4i^3 z + i^4 - 15 = (z+i)^4 - 16 = 0,
\]
ya que $(z+i)^4 = z^4 + 4iz^3 + 6i^2z^2 + 4i^3z + i^4$. Con la sustitución $w = z+i$ queda $w^4 = 16$, cuyas raíces son
\[
w_k = \sqrt[4]{16}\,e^{\,i\frac{2k\pi}{4}} = 2\,e^{\,i\frac{k\pi}{2}} \implies w_0=2,\ w_1=2i,\ w_2=-2,\ w_3=-2i.
\]
Deshaciendo $z = w - i$:
\[
z_0 = 2 - i, \qquad z_1 = i, \qquad z_2 = -2 - i, \qquad z_3 = -3i.
\]
\end{ejemplo}

\begin{truco}
\begin{itemize}
  \item \textbf{Cuidado con el cuadrante.} $\arctan(b/a)$ siempre devuelve un ángulo en $\left(-\tfrac{\pi}{2}, \tfrac{\pi}{2}\right)$. Antes de escribir el argumento, ubicá $z$ en el plano y sumá $\pi$ o $2\pi$ según el cuadrante. Un error de cuadrante arruina toda potencia o raíz.
  \item \textbf{Potencias grandes de $i$:} reducí el exponente módulo $4$. En general, para $z^{n}$ con $n$ grande, pasá a polar: $z^n = \rho^n e^{i n\theta}$ y reducí $n\theta$ módulo $2\pi$.
  \item \textbf{Dividir sin polar:} multiplicá numerador y denominador por el conjugado del denominador; el denominador se vuelve real $= \abs{w}^2$.
  \item \textbf{Contar raíces:} $w^n = z$ siempre da \emph{exactamente} $n$ raíces distintas (si $z\neq0$). Si te salen más, repetiste una.
  \item \textbf{Truco de conjugados:} en ecuaciones con $z$ y $\bar z$ mezclados, escribí $z = a+ib$ (o $z=\rho e^{i\theta}$, con $\bar z = \rho e^{-i\theta}$) y separá parte real e imaginaria para obtener un sistema en $\R$.
\end{itemize}
\end{truco}

\section{Vectores, producto interno, normas y ortogonalidad}

\begin{teoria}
\textbf{Vectores en $\K^n$.} Se trabaja sobre un cuerpo $\K$ que en este curso es $\R$ o $\C$. Un vector es una $n$-upla
\[
v = (v_1, v_2, \dots, v_n), \qquad v_i \in \K.
\]
\textbf{Operaciones} (componente a componente). Sean $u,v \in \K^n$, $\lambda \in \K$:
\begin{align*}
u = v &\iff u_i = v_i \ \ (1 \le i \le n), \\
u + v &= (u_1 + v_1, \dots, u_n + v_n), \\
\lambda\,v &= (\lambda v_1, \dots, \lambda v_n).
\end{align*}
La suma es ley de composición interna $f:\K^n\times\K^n \to \K^n$ y el producto por escalar una ley de composición externa $f:\K\times\K^n \to \K^n$. Cuando $v \in \R^n$ con $n=1,2,3$ el vector se llama \emph{geométrico}.
\end{teoria}

\begin{teoria}
\textbf{Producto interno.} Sean $u,v \in \K^n$. El producto interno (o interior) es
\[
\inner{u}{v} = \sum_{k=1}^{n} u_k\,v_k \qquad (\text{caso real, } \K = \R).
\]
En el caso hermítico ($\K = \C$) la definición usual, que garantiza $\inner{u}{u} \ge 0$, conjuga el segundo factor:
\[
\inner{u}{v} = \sum_{k=1}^{n} u_k\,\overline{v_k} \qquad (\text{producto hermítico}).
\]
\textbf{Propiedades} (caso real; en el hermítico la simetría pasa a ser $\inner{u}{v} = \overline{\inner{v}{u}}$):
\begin{itemize}
  \item \textbf{Simetría:} $\inner{u}{v} = \inner{v}{u}$.
  \item \textbf{Bilinealidad:} $\inner{u+w}{v} = \inner{u}{v} + \inner{w}{v}$; $\;\inner{\alpha u}{v} = \alpha\inner{u}{v}$ para $\alpha \in \K$ (análogo en el segundo argumento).
  \item \textbf{Positividad:} $\inner{u}{u} \ge 0$, con igualdad sólo si $u=0$.
\end{itemize}
\textbf{Desigualdad de Cauchy--Bunyakovsky--Schwarz:}
\[
\abs{\inner{u}{v}}^2 \le \inner{u}{u}\,\inner{v}{v} = \norm{u}_2^2\,\norm{v}_2^2 \qquad \forall\,u,v \in \K^n.
\]
\end{teoria}

\begin{teoria}
\textbf{Norma.} Una norma en $\K^n$ es una función $N:\K^n \to \R_{\ge 0}$ que cumple:
\begin{itemize}
  \item $N(v) \ge 0$ y $N(v)=0 \iff v=0$;
  \item $N(\lambda v) = \abs{\lambda}\,N(v)$ para todo $\lambda \in \K$;
  \item $N(u+v) \le N(u) + N(v)$ (desigualdad triangular).
\end{itemize}
\textbf{Norma euclidiana (usual, $L^2$):} inducida por el producto interno,
\[
\norm{u}_2 = \sqrt{\inner{u}{u}} = \sqrt{\sum_{i=1}^n \abs{u_i}^2}.
\]
Con esta norma $\K^n$ es un espacio normado.
\end{teoria}

\begin{teoria}
\textbf{Otras normas ($L^p$).} Sea $v = (v_1,\dots,v_n) \in \K^n$:
\[
\norm{v}_1 = \sum_{i=1}^n \abs{v_i} \quad (\text{norma 1}), \qquad
\norm{v}_p = \left(\sum_{i=1}^n \abs{v_i}^p\right)^{1/p} \quad (p > 1),
\]
\[
\norm{v}_\infty = \max\{\, \abs{v_i} : 1 \le i \le n \,\} \quad (\text{norma infinito}).
\]
La norma $2$ es el caso $p=2$; la norma infinito es el límite $p \to \infty$.
\end{teoria}

\begin{teoria}
\textbf{Ángulo, paralelismo y ortogonalidad.} El ángulo $\alpha_{uv}$ entre $u,v$ no nulos se define por
\[
\cos(\alpha_{uv}) = \frac{\inner{u}{v}}{\norm{u}_2\,\norm{v}_2}, \qquad 0 \le \alpha_{uv} \le \pi.
\]
(La desigualdad de Cauchy--Schwarz garantiza que el cociente está en $[-1,1]$.)
\begin{itemize}
  \item $u \parallel v$ (paralelos) si existe $k \in \K$ con $v = k\,u$.
  \item $u \perp v$ (perpendiculares/ortogonales) si son no nulos y $\inner{u}{v} = 0$.
\end{itemize}
\textbf{Proyección ortogonal.} La proyección de $u$ sobre $v$ ($v \neq 0$) es el vector
\[
\operatorname{proy}_v(u) = \frac{\inner{u}{v}}{\inner{v}{v}}\,v = \frac{\inner{u}{v}}{\norm{v}_2^2}\,v,
\]
y el escalar (proyección escalar) es $\dfrac{\inner{u}{v}}{\norm{v}_2}$. El residuo $u - \operatorname{proy}_v(u)$ es ortogonal a $v$.
\end{teoria}

\begin{receta}
\textbf{Analizar dos vectores $u,v \in \K^n$ (ángulo / ortogonalidad / proyección).}
\begin{enumerate}
  \item Calcular el producto interno $\inner{u}{v} = \sum u_i v_i$ (conjugando el segundo factor si $\K = \C$).
  \item Calcular las normas $\norm{u}_2 = \sqrt{\inner{u}{u}}$ y $\norm{v}_2 = \sqrt{\inner{v}{v}}$.
  \item \textbf{Ortogonalidad:} si $\inner{u}{v}=0$ (y ambos no nulos), son ortogonales, $\alpha = \tfrac{\pi}{2}$.
  \item \textbf{Ángulo:} $\alpha_{uv} = \arccos\!\left(\dfrac{\inner{u}{v}}{\norm{u}_2\norm{v}_2}\right)$.
  \item \textbf{Proyección de $u$ sobre $v$:} $\operatorname{proy}_v(u) = \dfrac{\inner{u}{v}}{\norm{v}_2^2}\,v$. Verificar que $u - \operatorname{proy}_v(u) \perp v$.
\end{enumerate}
\end{receta}

\begin{ejemplo}{Pizarrón Clase 2, Ej.\ 3 (adaptado)}
Sean $\vec a_1 = (1+i,\ 3-2i) \in \C^2$ y $\vec a_2 = (10+8i,\ 5+i\sqrt3) \in \C^2$.

\smallskip
\textbf{Parte real e imaginaria} (componente a componente):
\[
\operatorname{Re}(\vec a_1) = (1,\ 3), \quad \operatorname{Im}(\vec a_1) = (1,\ -2), \qquad
\operatorname{Re}(\vec a_2) = (10,\ 5), \quad \operatorname{Im}(\vec a_2) = (8,\ \sqrt3).
\]
\textbf{Suma:}
\[
\vec a_1 + \vec a_2 = \big(1+i+10+8i,\ \ 3-2i+5+i\sqrt3\big) = \big(11+9i,\ \ 8 + i(\sqrt3 - 2)\big).
\]
\textbf{Producto por escalar} con $\alpha = i\sqrt3$:
\[
\alpha\,\vec a_1 = i\sqrt3\,(1+i,\ 3-2i) = \big(i\sqrt3 - \sqrt3,\ \ 3i\sqrt3 + 2\sqrt3\big),
\]
donde se usó $i\sqrt3\cdot i = i^2\sqrt3 = -\sqrt3$ en la primera componente y $i\sqrt3\cdot(-2i) = -2i^2\sqrt3 = 2\sqrt3$ en la segunda.
\end{ejemplo}

\begin{truco}
\begin{itemize}
  \item \textbf{Ortogonalidad = producto interno cero.} Chequear ortogonalidad es sólo un producto interno; no hace falta calcular el ángulo entero.
  \item \textbf{En $\C^n$ conjugá el segundo factor.} Con el producto hermítico $\inner{u}{v} = \sum u_i \overline{v_i}$ se garantiza $\inner{u}{u} = \sum \abs{u_i}^2 \ge 0$ real. Si no conjugás, $\inner{u}{u}$ puede dar complejo y la norma no tiene sentido.
  \item \textbf{Normas ordenadas:} para todo $v$, $\norm{v}_\infty \le \norm{v}_2 \le \norm{v}_1$. Útil para acotar rápido.
  \item \textbf{Norma vía producto interno:} $\norm{u}_2^2 = \inner{u}{u}$; conviene trabajar con el cuadrado para evitar raíces hasta el final.
  \item \textbf{Proyección:} el numerador $\inner{u}{v}$ y el denominador $\norm{v}_2^2 = \inner{v}{v}$ usan el \emph{mismo} vector $v$ en el denominador, no $u$.
\end{itemize}
\end{truco}

\section{Matrices y determinantes}

\begin{teoria}
\textbf{Matriz y operaciones.} Una matriz $A \in \K^{n\times m}$ es un arreglo $A = (a_{ij})_{ij}$ con $a_{ij} \in \K$, $1\le i\le n$, $1\le j\le m$. Suma y producto por escalar son entrada a entrada:
\[
(A+B)_{ij} = a_{ij} + b_{ij}, \qquad (\lambda A)_{ij} = \lambda\,a_{ij}.
\]
\textbf{Producto de matrices.} Si $A \in \K^{n\times p}$ y $B \in \K^{p\times m}$, entonces $C = AB \in \K^{n\times m}$ con
\[
c_{ij} = \sum_{k=1}^{p} a_{ik}\,b_{kj} = \inner{F_i^A}{C_j^B},
\]
o sea la entrada $(i,j)$ es el producto de la fila $i$ de $A$ por la columna $j$ de $B$. \textbf{Cuidado:} el producto \emph{no} es conmutativo ($AB \neq BA$ en general); puede existir $AB$ y no $BA$. Es asociativo y distributivo. La potencia $A^m = A\cdots A$ ($m$ veces) requiere $A$ cuadrada.

\medskip
\textbf{Identidad.} $I_n = (\delta_{ij})$, con $1$ en la diagonal y $0$ fuera; cumple $AI_n = I_n A = A$.
\end{teoria}

\begin{teoria}
\textbf{Transpuesta, conjugada y adjunta (hermítica).}
\[
A^T = (a_{ji})_{ij} \quad (\text{transpuesta}), \qquad
\adjm{A} = \bar A^{\,T} = (\overline{a_{ji}})_{ij} \quad (\text{traspuesta conjugada / hermítica}).
\]
Propiedades de la transpuesta:
\[
(A^T)^T = A, \quad (A+B)^T = A^T + B^T, \quad (\lambda A)^T = \lambda A^T, \quad (AB)^T = B^T A^T.
\]
(análogas para $\adjm{A}$, con $(\lambda A)^{*} = \bar\lambda\,\adjm{A}$).
\end{teoria}

\begin{teoria}
\textbf{Clasificación de matrices cuadradas} ($A \in \K^{n\times n}$):
\begin{itemize}
  \item \textbf{Simétrica:} $A = A^T$ (es decir $a_{ij} = a_{ji}$).
  \item \textbf{Antisimétrica:} $A = -A^T$ ($a_{ij} = -a_{ji}$; en la diagonal $a_{ii}=0$).
  \item \textbf{Hermítica:} $A = \adjm{A}$ (autoadjunta; generaliza la simétrica al caso complejo).
  \item \textbf{Triangular superior:} $a_{ij}=0$ para $i>j$; \textbf{inferior:} $a_{ij}=0$ para $i<j$.
  \item \textbf{Diagonal:} $a_{ij}=0$ para $i\neq j$; \textbf{escalar:} $A = k\,I_n$.
  \item \textbf{Idempotente:} $A^2 = A$ (entonces $A^n = A$ para todo $n$).
  \item \textbf{Involutiva:} $A^2 = I$ (entonces $A^{2n}=I$, $A^{2n+1}=A$).
  \item \textbf{Nilpotente:} existe $n_0 \in \N$ con $A^{n} = 0$ para $n \ge n_0$ ($n_0$ = índice de nilpotencia).
  \item \textbf{Ortogonal:} $A A^T = A^T A = I_n$ (sus columnas forman base ortonormal; en el caso complejo, \textbf{unitaria:} $A\,\adjm{A} = I_n$).
\end{itemize}
\end{teoria}

\begin{teoria}
\textbf{Inversa.} $A \in \K^{n\times n}$ es invertible si existe $A^{-1}$ con $A\,A^{-1} = A^{-1}A = I_n$. La inversa, si existe, es única, y:
\[
(A^{-1})^{-1} = A, \qquad (\lambda A)^{-1} = \tfrac{1}{\lambda}A^{-1}\ (\lambda\neq0), \qquad (AB)^{-1} = B^{-1}A^{-1}.
\]
\end{teoria}

\begin{teoria}
\textbf{Determinante.} Es la única forma multilineal alternada $\det:\K^{n\times n}\to\K$ (lineal en cada columna, se anula con columnas repetidas) que cumple $\det(I_n)=1$. Notaciones: $\det(A) = \abs{A} = D(A)$.

\smallskip
\textbf{Casos chicos:}
\[
\det\begin{pmatrix} a & b \\ c & d \end{pmatrix} = ad - bc,
\]
y para $3\times 3$, la \emph{regla de Sarrus}:
\[
\det(A) = a_{11}a_{22}a_{33} + a_{12}a_{23}a_{31} + a_{13}a_{21}a_{32} - a_{13}a_{22}a_{31} - a_{11}a_{23}a_{32} - a_{12}a_{21}a_{33}.
\]

\smallskip
\textbf{Menor y cofactor.} El menor $M_{ij}$ es la submatriz $(n-1)\times(n-1)$ que resulta de borrar la fila $i$ y la columna $j$. El cofactor es
\[
C_{ij} = A_{ij} = (-1)^{i+j}\,\abs{M_{ij}}.
\]
\textbf{Regla de Laplace} (desarrollo por la fila $i$; análogo por columnas):
\[
\det(A) = \sum_{j=1}^{n} a_{ij}\,C_{ij}.
\]
\end{teoria}

\begin{teoria}
\textbf{Propiedades del determinante.} Sea $A \in \K^{n\times n}$:
\begin{itemize}
  \item Multilinealidad por columnas: $\det(\dots, \lambda A_i, \dots) = \lambda\det(\dots,A_i,\dots)$, y separa sumas.
  \item Intercambiar dos columnas (o filas) cambia el signo: $\det(\dots,A_i,\dots,A_j,\dots) = -\det(\dots,A_j,\dots,A_i,\dots)$.
  \item Sumar a una columna un múltiplo de otra \emph{no} cambia el determinante.
  \item Una columna nula, o dos columnas iguales o linealmente dependientes, dan $\det = 0$.
  \item $\det(A^T) = \det(A)$; $\;\det(AB) = \det(A)\det(B)$; $\;\det(kA) = k^n\det(A)$; $\;\det(A^{-1}) = \dfrac{1}{\det(A)}$.
  \item $A$ es invertible $\iff \det(A)\neq0$; \emph{singular} $\iff \det(A)=0$.
\end{itemize}
\textbf{Inversa por adjunta.} Si $\det(A)\neq0$,
\[
A^{-1} = \frac{1}{\det(A)}\,\operatorname{adj}(A), \qquad \operatorname{adj}(A) = (C_{ij})^T = (\text{matriz de cofactores transpuesta}).
\]
\end{teoria}

\begin{receta}
\textbf{Calcular $\det(A)$ y decidir singularidad.}
\begin{enumerate}
  \item Si $A$ es $2\times2$: $ad-bc$. Si es $3\times3$: Sarrus.
  \item Para $n\ge4$ (o con muchos ceros): desarrollo de Laplace por la fila/columna con \emph{más ceros}, para anular términos. Recordar el signo $(-1)^{i+j}$.
  \item Alternativa: triangular por operaciones de fila (sumar a una fila múltiplo de otra no cambia $\det$; intercambiar filas cambia el signo; escalar una fila por $\lambda$ multiplica $\det$ por $\lambda$). El determinante de una triangular es el producto de la diagonal.
  \item \textbf{Singular} si $\det(A)=0$; \textbf{invertible} si $\det(A)\neq0$. Si hay un parámetro $k$, plantear $\det(A)=0$ y resolver para $k$.
\end{enumerate}
\end{receta}

\begin{receta}
\textbf{Determinantes con columnas combinadas (usar propiedades, no calcular).}
\begin{enumerate}
  \item Escribir $\det$ de la matriz combinada en términos de las columnas originales $A_1,\dots,A_n$.
  \item Usar $C_i \to C_i + \alpha C_j$ (no cambia $\det$) para eliminar términos y dejar columnas originales.
  \item Factorizar escalares por columna: $\det(\dots,\lambda A_i,\dots) = \lambda\det(\dots)$.
  \item Contar intercambios de columnas: cada uno aporta un factor $(-1)$.
  \item Reducir a $\det(A_1,\dots,A_n) = \det(A)$ y sustituir el valor dado.
\end{enumerate}
\end{receta}

\begin{ejemplo}{Pizarrón Clase 3 / Guía II, Matrices Ej.\ 5}
Sea $A = (A_1, A_2, A_3) \in \R^{3\times3}$ con $\det(A) = -3$. Calcular $\det\!\left(\tfrac{3}{4}A^{-1}B^T\right)$, donde $B = (A_1 - 3A_3,\ A_3,\ A_1 - A_2)$.

\smallskip
\textbf{Separar factores externos.} Como $A^{-1}B^T \in \R^{3\times3}$:
\[
\det\!\left(\tfrac{3}{4}A^{-1}B^T\right) = \left(\tfrac{3}{4}\right)^3 \det(A^{-1})\det(B^T) = \tfrac{27}{64}\cdot\frac{1}{\det(A)}\cdot\det(B),
\]
usando $\det(kM)=k^n\det(M)$, $\det(MN)=\det M\det N$, $\det(A^{-1})=1/\det A$ y $\det(B^T)=\det(B)$.

\smallskip
\textbf{Calcular $\det(B)$ por operaciones de columna}, partiendo de $B = (A_1-3A_3,\ A_3,\ A_1-A_2)$:
\begin{align*}
C_1 \to C_1 + 3C_2:\quad &\det(B) = \det(A_1,\ A_3,\ A_1 - A_2), \\
C_3 \to C_3 - C_1:\quad &\det(B) = \det(A_1,\ A_3,\ -A_2) = -\det(A_1,\ A_3,\ A_2), \\
C_2 \leftrightarrow C_3:\quad &\det(B) = (-1)(-1)\det(A_1,\ A_2,\ A_3) = \det(A) = -3.
\end{align*}
\textbf{Resultado:}
\[
\det\!\left(\tfrac{3}{4}A^{-1}B^T\right) = \tfrac{27}{64}\cdot\frac{1}{-3}\cdot(-3) = \frac{27}{64}.
\]
\end{ejemplo}

\begin{truco}
\begin{itemize}
  \item \textbf{Elegí la fila/columna con más ceros} para Laplace: cada cero elimina un cofactor entero.
  \item \textbf{No calcules la inversa si sólo querés $\det(A^{-1})$:} es $1/\det(A)$. Idem, no multipliques matrices si sólo querés $\det(AB) = \det A\det B$.
  \item \textbf{El factor escalar es $k^n$, no $k$:} $\det(kA) = k^n\det(A)$, con $n$ el orden de la matriz. Error clásico.
  \item \textbf{Singular $\iff$ columnas linealmente dependientes} $\iff \det=0$. Para hallar el $k$ que vuelve singular, planteá $\det=0$.
  \item \textbf{Triangular es gratis:} el $\det$ de una triangular (o diagonal) es el producto de la diagonal. Si triangulás para calcular, recordá que intercambiar filas cambia el signo.
  \item \textbf{$(AB)^{-1} = B^{-1}A^{-1}$ y $(AB)^T = B^T A^T$:} el orden se \emph{invierte}. No es $A^{-1}B^{-1}$.
\end{itemize}
\end{truco}
